\section{Hybrid cloud}
\label{sec:Hybrid_cloud}

Hybrid cloud is one of the four cloud deployment models defined by the National Institute of Standards and Technology (NIST). According to Mell and Grance, a hybrid cloud is ``a composition of two or more distinct cloud infrastructures (private, community, or public) that remain unique entities but are bound together by standardized or proprietary technology that enables data and application portability'' \cite{Mell2011NIST}. Unlike public or private cloud deployments alone, hybrid cloud allows organizations to distribute workloads across multiple computing environments while maintaining interoperability between them. This deployment model enables organizations to leverage the scalability and elasticity of public cloud services while retaining sensitive applications and data within private cloud infrastructures to satisfy security, regulatory, or organizational requirements.

To facilitate interoperability among heterogeneous cloud infrastructures, the National Institute of Standards and Technology further proposed the Cloud Computing Reference Architecture (CCRA), which describes the major actors, functional components, and interactions within cloud computing systems \cite{Liu2011NISTRA}. The reference architecture provides a standardized view of how cloud providers, cloud consumers, cloud brokers, cloud auditors, and cloud carriers interact to support cloud services. Within hybrid cloud environments, these components cooperate to enable workload portability, resource orchestration, service management, and communication across multiple cloud platforms, forming the architectural foundation for modern hybrid cloud deployments.

The flexibility offered by hybrid cloud has led to its widespread adoption across both industry and academia. By combining public and private cloud resources, organizations can optimize infrastructure utilization, improve service availability, and dynamically allocate workloads according to business requirements while maintaining control over critical assets. Recent industry reports indicate that hybrid cloud has become the dominant deployment strategy for enterprise organizations, driven by the need to balance scalability, operational efficiency, regulatory compliance, and data governance \cite{Flexera2025StateCloud, HashiCorp2025StateCloud}. As cloud infrastructures continue to expand, organizations increasingly rely on hybrid cloud environments to support cloud-native applications, distributed services, and enterprise-scale software delivery.

Despite these advantages, hybrid cloud environments introduce significant operational complexity. Applications, services, and infrastructure components are distributed across heterogeneous platforms with different management interfaces, security policies, and deployment mechanisms. Maintaining consistent configurations, enforcing uniform security policies, and coordinating operational activities across multiple cloud environments become increasingly challenging as infrastructure complexity grows \cite{Liu2011NISTRA, HashiCorp2025StateCloud}. These characteristics have motivated the adoption of operational methodologies that integrate infrastructure management, software delivery, and continuous security throughout the software operational lifecycle. Among these methodologies, DevOps and its security-oriented evolution, DevSecOps, have become prominent approaches for improving automation and collaboration in cloud environments.